% Created 2017-10-03 Tue 16:10
% Intended LaTeX compiler: pdflatex
\documentclass[11pt]{article}
\usepackage[utf8]{inputenc}
\usepackage[T1]{fontenc}
\usepackage{graphicx}
\usepackage{grffile}
\usepackage{longtable}
\usepackage{wrapfig}
\usepackage{rotating}
\usepackage[normalem]{ulem}
\usepackage{amsmath}
\usepackage{textcomp}
\usepackage{amssymb}
\usepackage{capt-of}
\usepackage{hyperref}
\usepackage[margin=0.80in]{geometry}
\usepackage[backend=biber, style=ieee]{biblatex}
\usepackage{float}
\usepackage[super,negative]{nth}
\usepackage[capitalise]{cleveref}
\usepackage{pst-node,transparent,ragged2e}
\addbibresource{/home/nuk3/course/training/csir/novellasers/bibliography/bibliography.bib}
\DeclareFieldFormat[inproceedings]{citetitle}{\textit{#1}}
\DeclareFieldFormat[inproceedings]{title}{\textit{#1}}
\DeclareFieldFormat[misc]{citetitle}{#1}
\DeclareFieldFormat[misc]{title}{#1}
\renewcommand*{\bibpagespunct}{%
\ifentrytype{inproceedings}
{\addspace}
{\addcomma\space}}
\AtEveryCitekey{\ifuseauthor{}{\clearname{author}}}
\AtEveryBibitem{\ifuseauthor{}{\clearname{author}}}
\SpecialCoor
\author{Nyameko Lisa}
\date{22 September 2017}
\title{IPSP044 - Essential Trademark Law\\\medskip
\large Assignment 02 - 869602}
\hypersetup{
 pdfauthor={Nyameko Lisa},
 pdftitle={IPSP044 - Essential Trademark Law},
 pdfkeywords={},
 pdfsubject={},
 pdfcreator={Emacs 26.0.50 (Org mode 9.0.9)},
 pdflang={English}}
\begin{document}

\maketitle
\addvspace{110pt}
\centering{
\pnode(0.5\textwidth,-0.5\textheight){thisCenter}
\rput(thisCenter){%\transparent{0.25}
\includegraphics[width=2.7in]{/home/nuk3/course/llb/wipo-unisa/UNISACoatofArms.eps}}}

\justifying
\addvspace{110pt}
\section*{Declaration}
\label{sec:org3ef4e75}
I know that plagiarism is to use someone else’s work and pass it off as my own.\\
I know that plagiarism is wrong.\\
I confirm that this assignment is my own work.\\
I have acknowledged in the bibliography accompanying the assignment all the sources that I have used.\\
I have not directly copied without acknowledgement anything from the Internet or from any other source.\\
I have indicated every quotation and citation in a footnote or bracket linked to that quotation.\\
I have not allowed anyone else to copy my work and to pass it off as their own work.\\
I understand that if any unacknowledged copying whatsoever appears in my assignment I will receive zero per cent for the assignment.\\
I am aware of the UNISA policy on plagiarism and understand that disciplinary proceedings can be instituted against me by UNISA if I contravene this policy.\\
I indicate my understanding and acceptance of this declaration by
entering my name here under:
\begin{itemize}
\item Name: \textbf{Nyameko Lisa} (Student Number: \textbf{7874-909-3})
\end{itemize}

\subsection*{NOTE}
\label{sec:org953d6bc}
Please note that footnotes will be denoted as \footnote{This is a footnote.} and will
appear at the bottom of the page.\\
References will be denoted by \cite{rsa93_tm_act} and will appear at the end of the document.
\newpage
\section{Advise F whether there are any grounds on which he can oppose V's application for the registration of the ELEKTROLIN mark in respect of electrical apparatus and instruments. [10]}
\label{sec:org7adb15a}
As per the definitions of [section 2(1)]\cite{rsa93_tm_act} and given that the Republic is a Paris Convention Member State, [article 15(1)]\cite{wto17_trips}, [article 1(2)]\cite{wipo83_paris_conve_protect_ip}, it follows that \emph{`LECTROLON'} and \emph{`ELEKTROLIN'} will be understood to be \textbf{marks}, in that they are names or signs capable of being represented graphically. Moreover they constitute \textbf{trade marks}, as they have been proposed by F and V respectively, in relation to distinguishing their trade or use from other goods or services connected in the course of trade regarding \emph{`electronic equipment'} or \emph{`electrical apparatus and equipment'}.\\

As per [section 2(4)]\cite{rsa93_tm_act}, \emph{`use'} or \emph{`proposed use'} shall be construed as the use of a visual representation of the mark [section 2(2)(a)]\cite{rsa93_tm_act}, as pertaining to physical form or other relation in terms of goods[section 2(3)(a)]\cite{rsa93_tm_act}, or as pertaining to the use thereof in any relation to the performance of any services\footnote{Also referred to as service marks, [article 1(2)]\cite{wipo83_paris_conve_protect_ip}} [section 2(3)(b)]\cite{rsa93_tm_act}.\\

It should be noted that LECTROLON does \textbf{NOT} constitute a \textbf{registered trade mark}, [section 2(1)]\cite{rsa93_tm_act}, and therefore as per [section 33]\cite{rsa93_tm_act} F is \textbf{unable} to bring an action for infringement against V, should V have already started to trade using his proposed trademark, ELEKTROLIN.\\

As per [section 21]\cite{rsa93_tm_act} and [article 15(5)]\cite{wto17_trips}, F has \textbf{3 months} to oppose V's application for the registration of the ELEKTROLIN mark in respect of electrical apparatus and instruments, from the date of the advertisement of the acceptance of the application, [section 17]\cite{rsa93_tm_act}. The basis of F's opposition will be the conflict of V's trademark with his own unregistered mark, it's general lack of inherent eligibility for trade mark registration and the exploitation of the goodwill that the LECTROLON mark has acquired through sales, advertising and promotion.\\

In establishing inherent ineligibility of V's mark, registration may be denied
on a number of grounds, however Paris Union Member States [article
1(1)]\cite{wipo83_paris_conve_protect_ip}, may not derogate from the provisions of
the Paris Convention, [article 15(2)]\cite{wto17_trips}. Furthermore, F must
demonstrate that V's application to register his trademark does \uline{not} amount to \textbf{honest concurrent use}, [section 14(1)]\cite{rsa93_tm_act}.\\

As per the provisions of [section 10(1)]\cite{rsa93_tm_act}, F could argue that
V's mark is an unregistrable  trade mark, in that ELEKTROLIN is incapable of
generally distinguishing the goods and services it provides, from those of
LECTROLON, [section 9(1) and section 10(2)(a)]\cite{rsa93_tm_act}. Moreover, F could claim \emph{`proprietorship'} of the ELEKTROLIN mark, and deny V \emph{bona fide} claim to proprietorship, [section 10(3)]\cite{rsa93_tm_act}.\\

As per the provisions of [section 10(6)]\cite{rsa93_tm_act}, F could argue that the ELEKTROLIN mark is unregistrable, due to the likelihood to cause deception or confusion, given that it constitutes a reproduction or imitation of his own mark, LECTROLON, which is entitled to protection under the Paris Convention, \cite{wipo83_paris_conve_protect_ip}. Where the requirement for protection of well-known marks under the Paris Convention will be satisfied provided that F is himself either a national, domiciled in, or has real and effective industrial or commercial establishment in a Paris Convention member state [section 35(1)]\cite{rsa93_tm_act}. Additionally, LECTROLON has garnered a reputation though sales, marketing and promotion in the relevant sector, it will indeed be considered to be \emph{``well-known''}, [section 35(1A)]\cite{rsa93_tm_act}.\\

Furthermore, it could be argued that since P has not granted any consent, if
registration of the ELEKTROLIN mark is permitted, it is likely to cause
deception or confusion with the registered trade mark LECTROLON, [section
10(12-13)]\cite{rsa93_tm_act}. Where the notion of \emph{`reasonable probability or
"likelihood" to cause deception or confusion'}, was extensively argued in
\citetitle{harms01_cowbell_v_ics}, \cite{harms01_cowbell_v_ics}, in accordance with
our case law and referencing similar cases, it was said that likelihood of
confusion must \emph{`be appreciated globally'} and that:
\begin{quote}
\emph{``\ldots{}global appreciation of the visual, aural or conceptual similarity of the marks in question, must be based on the overall impression given by the marks,
bearing in mind, in particular, their distinctive and dominant components.''}
\end{quote}


Moreover, as per the provisions of [article 2(1)]\cite{wipo96_model_provi_unfair_comp} and [article 10\(^{\text{bis}}\)(3)(3)]\cite{wipo83_paris_conve_protect_ip}, F could argue that use and registration of V's mark ELEKTROLIN, would amount to \emph{unfair competition}, in that its registration or use, during the course of trade is likely to cause confusion with the, \emph{unregistered} trademark, LECTROLON, [article 2(2)(i)]\cite{wipo96_model_provi_unfair_comp}, or to mislead the public [article 4(1)]\cite{wipo96_model_provi_unfair_comp} as to the nature of those goods, through exploitation of the goodwill enjoyed by F and his unregistered trademark LECTROLON. Thus amounting to \textbf{unfair competition} through the \emph{dilution}\footnote{\emph{``Lessening of the distinctive character or advertising value of a trademark''.}} of the unregistered trademark's goodwill and reputation, [article 3(2)(a)(i)]\cite{wipo96_model_provi_unfair_comp}.

\printbibliography
\end{document}